\section*{Chapitre \ref{ch:s1:news-and-events} --- Actualités et événements}

\begin{solution}{exo:s1:news-and-events:1}
$e^{-0.1/0.2} = e^{-0.5} = 0.607$ à 100 millisecondes~; $e^{-1/0.2} = e^{-5} = 0.0067$ à une seconde, soit moins d'un pour cent.
\end{solution}

\begin{solution}{exo:s1:news-and-events:2}
$1/(1 + 0.03 \times 39) = 0.46$, donc la part immédiate vaut $0.5 + 0.5 \times 0.46 = 0.73$ et 27~\% relèvent de la dérive. Pour l'attention médiane, $0.5 + 0.5 \times 0.42
= 0.71$.
\end{solution}

\begin{solution}{exo:s1:news-and-events:3}
Tant que l'action est suspendue, personne ne peut la négocier~; le mouvement immédiat se produit à la réouverture, lors de l'enchère d'ouverture, où chaque ordre rencontre un prix unique. Il n'existe
aucune fenêtre pendant laquelle un trader rapide puisse négocier à l'ancien prix.
\end{solution}

\begin{solution}{exo:s1:news-and-events:4}
La corrélation entre le ton et la lecture vaut $1/\sqrt{1 + 1.5^2} = 0.555$, donc les signes concordent avec la probabilité $\tfrac12 + \arcsin(0.555)/\pi = 0.687$~;
la simulation mesure 68.9~\%.
\end{solution}

\begin{solution}{exo:s1:news-and-events:5}
Au bout d'une seconde, la part immédiate a presque entièrement disparu ($e^{-5}$ en subsiste) et la dérive n'a pas encore commencé~: elle se manifeste dans les jours qui suivent.
Entre une seconde et la clôture, rien ne bouge dans le modèle.
\end{solution}

\begin{solution}{exo:s1:news-and-events:6}
Chaque moitié contient environ la moitié des dépêches, si bien que son P\&L quotidien porte en moyenne sur deux fois moins de positions indépendantes et que son bruit est plus élevé. L'ensemble du livre
se diversifie sur les deux moitiés~; sa dérive par dépêche (25.4 points de base) se situe entre les deux.
\end{solution}

\begin{solution}{exo:s1:news-and-events:7}
Le ratio de Sharpe monte à 7.6 et le rendement à 17.4~\% par an, avec 45.7 points de base de dérive par dépêche avant coûts, au lieu de 25.4. La différence tient à la
valeur du modèle de texte. Le trader conserve $25.4/45.7 = 56~\%$ de la dérive, soit la corrélation entre la lecture et le ton (0.555)~: plus que le
$2 \times 0.689 - 1 = 38~\%$ que laisse supposer le seul taux de réussite, car les mouvements les plus forts sont lus correctement plus souvent.
\end{solution}

\begin{solution}{exo:s1:news-and-events:8}
L'heure de publication d'un site n'est pas celle où l'information a atteint le marché~: les fils d'agence et les \omterm{def:s1:news-and-events:mrn}{dépêches lisibles par machine} la transportent plus tôt, et les machines la négocient
en moins d'une seconde. Un backtest calé sur la minute de publication peut négocier à des cours qui ont déjà bougé, ou, si l'heure du site est antérieure à sa mise à jour, avant que la
nouvelle n'existe. Utiliser les horodatages propres au flux, tels que reçus, puis les cours qui leur sont postérieurs.
\end{solution}

\begin{solution}{pb:s1:news-and-events:1}
\begin{enumerate}
\item Un flux de nouvelles structuré destiné aux programmes (horodatage, titres, catégorie, scores)~; l'intervalle allant de la dépêche, augmenté de la latence, jusqu'à la fin de la
période de détention.
\item Associer la dépêche à ses titres~; noter son sens et son intensité~; décider du moment d'agir.
\item Un pessimisme élevé annonce une pression baissière sur les cours, suivie d'un retour aux fondamentaux, et un pessimisme extrême un volume élevé.
\item La part de mots négatifs annonce des résultats faibles~; les cours sous-réagissent brièvement, surtout pour les dépêches portant sur les fondamentaux.
\item La réaction peut être préparée, ce qui récompense la rapidité~; la distraction due aux autres annonces est prévisible.
\item Réactions immédiates plus faibles et dérive plus forte lorsque de nombreuses sociétés annoncent le même jour~; alphas substantiels.
\item Une \omterm{def:s1:news-and-events:halt}{suspension de cotation} en attente d'une nouvelle, ou après celle-ci, ou après un fort mouvement~; T1~: nouvelle en attente, T2~: nouvelle publiée, pause de volatilité LUDP.
\item La rapidité durant la première seconde ne vaut rien~; le mouvement se produit à l'enchère de réouverture.
\item 1\,000 actions pendant huit ans~; ton avec un écart-type de 2~\%~; attention $1/(1+0.03(k-1))$~; une part immédiate comprise entre la moitié et la totalité, avec une constante de temps
de 0.2 secondes~; le reste sur cinq jours~; \omterm{def:s1:news-and-events:halt}{suspensions de cotation} au-delà de 5~\%.
\item 82\,351 dépêches, 40.8 par jour, dont 1.18~\% suivies d'une \omterm{def:s1:news-and-events:halt}{suspension de cotation}.
\item 96.7~\%, 70.1~\%, 29.2~\% et 28.7~\%.
\item En dessous d'une seconde, les machines prennent la part immédiate~; au-delà, seule la dérive subsiste.
\item \textbf{Résultat clé.} La part du mouvement encore à venir vaut 96.7~\% à 1 milliseconde, 93.7~\% à 10, 70.1~\% à 100, 29.2~\% à une seconde et 28.7~\% à
la clôture suivante~; un trader à horizon journalier, doté d'un modèle de texte imparfait, capte 25.4 points de base par dépêche, avec un ratio de Sharpe de 3.24 après coûts (1.80 les jours calmes,
2.74 les jours chargés).
\item L'attention se répartit sur davantage de dépêches, si bien que moins de chaque mouvement se produit d'un coup~: 29.5 contre 20.8 points de base de dérive par dépêche.
\item 44~\% de la dérive (il en conserve 25.4 sur 45.7 points de base), et le ratio de Sharpe passe de 7.6 avec une lecture parfaite à 3.24.
\item Le modèle introduit proprement la sous-réaction et rend indépendantes les erreurs de lecture~; la dérive réelle se mesure en points de base par dépêche, et les coûts en absorbent une grande partie
de celle-ci.
\item Horodatages du flux tels que reçus, cours postérieurs à ceux-ci, scores de ton issus de données point-in-time et coûts appliqués à chaque dépêche.
\item Le sentiment des titres lorsqu'il est négocié dans la première seconde~; la réouverture consécutive à une suspension, pour ses données d'enchère.
\item Le chapitre~7 a négocié la dérive après les résultats dans une \omterm{def:s1:earnings:event}{étude d'événement}~; ce chapitre mesure la part du mouvement de chaque dépêche qui reste à venir pour chaque vitesse.
\item Le trader rapide achète la part immédiate du mouvement~; les traders plus lents se partagent ce que l'attention laisse derrière elle.
\end{enumerate}
\end{solution}

\begin{solution}{iq:s1:news-and-events:1}
Associer chaque dépêche à ses titres et ne garder que les dépêches pertinentes et inédites~; noter son ton à l'aide d'un dictionnaire ou d'un modèle entraîné~; fixer la fenêtre de réaction et
l'horizon~; tester l'IC du signal par horizon sur des données point-in-time~; dimensionner la position après coûts.
\end{solution}

\begin{solution}{iq:s1:news-and-events:2}
L'heure à laquelle le flux a livré chaque dépêche à la société, à comparer à l'heure indiquée par l'article ou le fournisseur~; l'heure des cours utilisés pour l'entrée en position~; et si
les scores de ton ou les étiquettes ont été révisés a posteriori.
\end{solution}

\begin{solution}{iq:s1:news-and-events:3}
Lire les dépêches dès leur publication, estimer le cours de réouverture à partir du cours indicatif et du déséquilibre de l'enchère, décider si la position doit être conservée,
la réduire ou la couvrir au moyen d'instruments liés, et placer des ordres dans l'enchère de réouverture plutôt que de courir après les premières transactions qui suivent.
\end{solution}

\begin{solution}{iq:s1:news-and-events:4}
Un flux direct décodé dans un processus voisin du moteur de négociation, des correspondances précalculées entre entités et titres, un modèle de scoring assez léger pour tourner en
quelques microsecondes, des ordres préparés à l'avance avec des contrôles de risque effectués en amont, une colocation avec la bourse, et des horodatages à chaque étape pour mesurer la latence.
\end{solution}

\begin{solution}{iq:s1:news-and-events:5}
Les investisseurs ont une attention limitée~: avec de nombreuses annonces le même jour, chacune en reçoit moins, si bien que les cours réagissent moins immédiatement et dérivent davantage ensuite, comme
l'ont établi Hirshleifer, Lim et Teoh.
\end{solution}

\begin{solution}{iq:s1:news-and-events:6}
À latence $L$, la part immédiate encore à venir vaut $a e^{-L/\tau}$ et la dérive $1 - a$ n'est pas touchée, donc la part encore à venir vaut $a e^{-L/\tau} + 1 - a$.
La moitié de la part immédiate est épuisée lorsque $e^{-L/\tau} = \tfrac12$, soit $L = \tau \ln 2$~: 0.14 secondes pour $\tau = 0.2$.
\end{solution}
